\chapter*{About this manual}
\addcontentsline{toc}{chapter}{About this manual}
\markboth{About this manual}{About this manual}

This manual is the architectural specification of \IS, a prototype 64-bit instruction set architecture,
and of the platform, firmware interface and tools built around it. It is written for people who write
\IS software (compilers, operating systems, firmware, applications in assembly language) and for people
who implement the architecture in hardware or in simulators.

\section*{Structure}
\begin{description}[leftmargin=1em,style=nextline]
\item[Part~\ref{part:arch}, Architecture] The programmer's model: architectural state, the storage
  tiers and the Transfer Unit, capabilities and revocation domains, inter-core communication,
  software-managed caching and code paging, privilege modes, traps and events, and the memory model with
  its soundness argument.
\item[Part~\ref{part:isa}, Instruction Set] Instruction formats, the opcode map, the pseudocode
  conventions, and one entry per instruction giving its encoding, assembler syntax, operation and
  exceptions.
\item[Part~\ref{part:platform}, Platform and System Software] The system memory map, the device space,
  the platform table, the console, the block device, the timer and the framebuffer; the firmware
  interface (SFI) and boot protocol; the StrataOS operating system, a microkernel with preemptive
  multitasking, its servers, shell and window server.
\item[Part~\ref{part:tools}, Tools and Implementation] The assembler and image format, the simulators
  and the JIT, the conformance suite, the emulator in a web browser, annotated example programs,
  Strata Forth, and the prototype and FPGA implementations.
\item[Appendices] Calling convention, trap causes, CSRs, encoding tables, an alphabetical instruction
  index, a glossary, and known limitations.
\end{description}

\section*{Normative sources}
The architecture is defined by its reference emulator, \texttt{emu/emu.c}, which executes the
sequential semantics directly: every transfer completes at the instant it issues. Where this manual and
the emulator disagree, the emulator is authoritative and the manual is in error. The SystemVerilog
implementation in \texttt{rtl/} is one implementation of the architecture and is not normative; where
it differs from this manual the difference is a defect of that implementation
(Section~\ref{sec:rtl-deltas}). The conformance suite (Section~\ref{sec:conformance}) checks both
against the rules of this manual.

\section*{Status}
This is the \emph{prototype edition, version 0.1}, the first public release of the manual. It
describes the instruction set as the reference emulator, the RTL implementation and the conformance
suite agree on it, together with the prototype platform, StrataOS 0.1 and the tools. The ISA is not
frozen: encodings, CSR numbers and platform layouts may change between editions. Resource and
performance figures marked \estimate are projections.

\section*{Typographic conventions}
\begin{tabularx}{\linewidth}{@{}l X@{}}
\toprule
\textbf{Convention} & \textbf{Meaning}\\\midrule
\insn{tget}, \reg{a0}, \reg{c3} & Instruction mnemonics, register names and assembler text.\\
\csr{tvec} & Control and status register names.\\
\cause{capop} & Trap cause names.\\
\perm{R}, \perm{W}, \perm{X}, \perm{A}, \perm{D}, \perm{V} & Capability permissions.\\
\field{rd}, \field{imm16} & Fields of an instruction encoding.\\
\hex{4000} & Hexadecimal numbers. Plain numbers are decimal.\\
\texttt{x<h:l>} & Bits $h$ down to $l$ of value \texttt{x}; bit 0 is least significant.\\
KiB, MiB & $2^{10}$ and $2^{20}$ bytes.\\
\bottomrule
\end{tabularx}

Byte order is little-endian throughout: in memory, in the Local Stores, in images and in device
data structures. Sizes: a \term{byte} is 8 bits, a \term{halfword} 16, a \term{word} 32, a
\term{doubleword} 64.

\section*{Normative language}
\term{Must} and \term{must not} state requirements on software or implementations. \term{Is} and
\term{does} describe architectural behaviour. \term{Should} is a recommendation. Behaviour described as
\term{timing-visible} may differ between implementations and between runs.
